% !TeX root = ../papers/nc-p3.tex
\section{PI正則窓における頂点単純との拡張による中心有限性の判定}
\label{sec:pi-vertex-extension-criterion}
\begin{theorem}\label{thm:pi-vertex-extension-criterion}
任意の標準正則窓の折り畳み $B$ がPIとする。次は同値である。
\begin{enumerate}
\item $B$ は左右Noetherである。
\item $B$ は有限生成中心上有限である。
\item 各非頂点単純 $S$ と各頂点単純 $k_i$ について、次数を忘れた
$\operatorname{Ext}^1_B(S,k_i)$ と $\operatorname{Ext}^1_B(k_i,S)$ が双方零である。
\end{enumerate}
従って行列恒等式のサイズを固定すれば、第三条件を単純表現の有限個の行列階数条件で調べられる。
第三条件を満たす標準有限Frobenius候補で、指定サイズの恒等式を満たすことが分かっているものには、
前節の有限検査が正則性自体の必要十分条件になる。十五次元閉包の有限四辺では564次検査で足りる。
\end{theorem}
\begin{proof}[単純の分離とtrace環の加群]
第一条件から第二条件は前定理、逆は有限加群性による。第二条件の下では非頂点単純に非零作用する
正次数中心元がある。実際中心の正次数イデアルで割った環は有限次元次数付きで、頂点単純だけを持つ。
中心指標が異なれば全Extが消える。以下第三条件を仮定し、既約性も半素性も仮定しない。

非頂点単純 $S$ について $\operatorname{Tor}^B_j(S,\bigoplus k_i)=0$ が全 $j$ で成立する。
ここでテンソルの右側には左頂点単純を取る。双対は右加群の $\operatorname{Ext}^j_B(S,\bigoplus k_i)$ である。
両側Koszul分解の4次元Calabi--Yau双対により、有限次元加群の第 $j$ Extと逆順の第 $4-j$ Extは双対である。
これは両側分解の項に二つの加群を入れ、有限次元行列間のtrace対合で複体を双対化してもよい。
従って第0,1,3,4群は零である。全頂点を足したKoszul複体の鎖群の次元は
$\dim S$ の $1,4,6,4,1$ 倍で、Euler数が零なので第2群も零となる。左でも同じである。

$R=B/P$ を無限次元の斉次素商、$Q_R$ をその中心単純商とする。$T$ を $R$ に $Q_R$ の全ての被約traceを
添加した環、$C$ をそのtraceが生成する可換環とする。前節の行列不変式の議論で $C$ は有限型で $T$ は
$C$ 上有限、また斉次な非零 $x\in Z(R)$ で $xT\subset R$ となるものがある。
$C_0=k$ である。次数零の語は頂点冪等元の一次結合であり、そのtraceは定数である。
任意の $\phi\in\operatorname{Spec} C$ の閉点で、$T_\phi$ の単純は $R$ に制限しても単純である。
実際 $T_\phi$ は $R$ の像で生成されている。

$\phi$ が錐の頂点でなければ $T_\phi$ の全ての単純は $B$ の非頂点単純であることを示す。
$C$ の一般点から $\phi$ へ特殊化する離散付値環を取り、必要なら剰余体を拡大し、一般の
$Q_R$ を分裂させる有限拡大と完備化を行う。例えば $\phi$ の爆上げの例外因子の一般点の付値を
使える。あるいは一般の中心単純開集合に一般点を持ち $\phi$ を通る曲線を取ってもよい。
$T$ は中心 $C$ 上有限なので一般の絶対単純表現には $T$ 安定な格子がある。自由格子に付値環上有限な
$T$ の像を作用させた和を取ればよい。その特殊繊維の半単純化には $T_\phi$ の全ての単純型が出る。
実際この表現のtraceは $T_\phi$ の特殊化traceである。もし単純型を欠けば、その型のみを台とする非零冪等元を
半単純商から持ち上げるとその全ての正冪のtraceが零になる。Newton公式と被約Cayley--Hamilton恒等式は
特殊化後にも成り立つから、この冪等元が冪零となり矛盾する。体拡大後にも同じである。
全単純が頂点単純であることもない。そうなら全ての正次数の語のtraceが零で $\phi$ は頂点となる。

両種が共存したと仮定する。第三条件と組成列による帰納から特殊格子表現は頂点型だけを持つ部分と
非頂点型だけを持つ部分の直和に分かれ、その二部分間の両向きの第一Extは零である。
この分解は完備付値環に持ち上がる。実際 $t^l$ まで対角化した矢行列の次の非対角項は、関係式の
非対角項を取ると二つの特殊ブロックの間の拡張のコサイクルとなる。零Extゆえ $1+t^l H$ の
頂点基底変更で非対角項を除ける。完備極限で一般表現を分解し、一般の絶対単純性に反する。
体拡大を使ったときも第三条件は保存される。固定された次元の単純性と非零Extの階数条件は有限型の
構成可能条件なので、代数閉拡大で反例があれば元の代数閉体にも反例がある。

従って任意の有限生成次数付き右 $T$ 加群 $N$ は $B$ に制限しても有限生成射影による有限分解を持つ。
実際 $N$ は $C$ 上有限で、頂点複体による $\operatorname{Tor}^B_j(N,\bigoplus k_i)$ も $C$ 上有限である。
非頂点閉点 $\phi$ の全ての厚み $N/\mathfrak m_\phi^l N$ の単純因子は上記の非頂点型なので、
全Torは消える。逆極限で完備化した複体も完全である。有限次元ベクトル空間の逆系では核もMittag--Lefflerであり、
また有限 $C$ 加群の完備化は完全である。よって全てのTorの台は頂点に限られ有限次元となる。
下に有界な次数付き加群の最小分解を逐次作れば、有限なTorが各段の生成元を与え、大域次元4で止まる。
左加群も同様である。
\end{proof}

\begin{proof}[任意の片側イデアルに現れる加群]
上の任意の $N$ について $d=\dim_C\operatorname{Supp}N$（すなわち台のKrull次元）と書けば、
$j<4-d$ で $\underline{\operatorname{Ext}}_B^j(N,B)=0$ である。
$d=0$ は有限長の場合である。$d>0$ では頂点だけに台を持つ最大部分を除いてよい。
これは有限長で、第4未満のExtの消滅を使える。残りには正次数の斉次な $C$ の非零因子を選べ、
余核は台の次元が小さい有限 $T$ 加群となる。長完全列で必要なExtは正次数移動と同型になるが、
前段の有限射影分解により下に有界なので零である。

各非零な斉次部分加群 $M\subset e_iB$ の中に、以上の形の非零 $N$ が存在することを示す。
任意のアフィン次数付きPI商で有限個の斉次素イデアルのある積が零である。
実際Jacobson根基は冪零で、その半素商の極小素は定理\ref{thm:finite-semiprime-specialization}冒頭の
有界行列表示の議論で有限個である。それらの積の適当な冪を取る。
右からこの素イデアル列を順に掛け、零となる直前の非零部分加群を取ると、一つの斉次素商 $R$ 上で作用する
非零部分加群 $M'$ を得る。$R$ が有限次元なら有限長部分加群が生じてAS消滅に反する。
前段の中心分母 $x$ の $M'$ 上の核が非零なら、その核は $R/xR$ 上の加群なので、そこで同じ素商選択を繰り返す。
この操作は有限回で終わる。$R\subset T$ と $xT\subset R$ より $R$ のGK次元は $C$ の次元 $e$ に等しく、
各斉次片の次元は $O(n^{e-1})$ である。斉次非零因子 $x$ で割った累積次元は末尾の $\deg x$ 個の斉次片の
和であり、従って次の商のGK次元は少なくとも1下がる。分母が次数零のときは初めから単元である。
最後は $x$ が単射となる。そこで $M'$ をその $x$ 局所化に埋め、斉次な非零元 $m\in M'$ を選ぶと
\[
                       0\ne N=mxT\subset M'
\]
となる。$T\subset R[1/x]$ と $xT\subset R$ を使った。この $N$ は一生成の次数付き右 $T$ 加群である。
包含が非零なのでExt消滅から $d\ge4$。一方 $M$ 自体の成長が高々4なので $d=4$ である。
有限射影分解により $H_N(t)=p(t)/(1-t)^4$ の $p(1)$ は正整数である。

頂点射影間の非零な斉次経路元の左乗法を取る。核も非零なら像と核の双方にこのような $N$ が入り、
二つのHilbert関数の和が始点の頂点射影のもの以下になる。最高係数の整数は両方正でその和が1以下にはならない。
よって左乗法は単射、全ての非零合成が非零であり $B$ は素である。$R=B$ に戻るとそのtrace拡大 $T$ 自体が
前段により左右完全である。補題\ref{lem:cy-finite-birational-transfer}から $T=B$、従って中心上有限である。
最後に全単純の次元は行列恒等式で有界であり、単純性・一次拡張は有限箙表示の連立式で調べられる。
\end{proof}
